\chapter{The software development lifecycle}
\label{ch:sdlc}

Blendity is written with an AI under a person's direction and review. This chapter is the lifecycle as it is
practised, with every artifact named by its path, so that someone can repeat it or audit it.

\section{The roles}
\begin{itemize}
  \item \textbf{The user} (Emiliano Manalo) decides what Blendity is for, which features it gets and in what order,
    approves each specification and decision record before code is written, reviews every pull request on GitHub and
    merges it, and decides when a release is cut. Nothing lands without that review (C-05, \adr{0006}).
  \item \textbf{The main session} (Claude through Claude Code \cite{ClaudeCode2026}) plans: it reads
    \code{docs/requirements.md}, \code{docs/ARCHITECTURE.md} and the decision records, writes the task spec and any
    ADR, implements the hard, cross-cutting or geometric work itself, runs the suites and the audits, writes the
    documentation and opens the pull request.
  \item \textbf{The \code{implementer} agent} implements a routine, well-specified task from its spec, builds and runs
    the unit tests; it makes no architectural decisions, adds no dependency, never commits.
  \item \textbf{The \code{test-engineer} agent} writes independent tests from the spec's acceptance criteria, edge cases
    first; it edits only the test folders and reports bugs as failing tests.
  \item \textbf{The \code{code-reviewer} agent} reviews the diff against the spec, the ADRs and \code{CLAUDE.md} and
    reports ranked findings with concrete failure scenarios and a verdict; it never edits files. Its findings are fixed
    or answered in the task's record. Agents run one after another, never fanned out.
\end{itemize}

\section{The lifecycle}
\begin{figure}[H]
  \centering
  \resizebox{\textwidth}{!}{% The agentic software development lifecycle: one task from request to release, with every artifact named by path.
\begin{tikzpicture}[node distance=4mm and 6mm]
  \node[person] (ask) {The user asks\\ \footnotesize a round of requests};
  \node[agent, right=of ask] (plan) {Plan (the main session)\\ \footnotesize reads \code{requirements.md},\\ \footnotesize \code{ARCHITECTURE.md}, the ADRs};
  \node[artifact, right=of plan] (spec) {\code{docs/tasks/NNNN-slug.md}\\ status: draft $\to$ approved\\ \code{docs/decisions/NNNN-*.md} (an ADR)\\ when the choice is architectural};
  \node[gate, right=of spec] (g1) {The human\\ approves the\\ spec and ADR};
  \node[agent, below=8mm of g1] (impl) {Implement\\ \footnotesize \code{implementer} (routine) or\\ \footnotesize the main session (hard)};
  \node[agent, left=of impl] (test) {Test\\ \footnotesize \code{test-engineer}: tests from the\\ \footnotesize acceptance criteria};
  \node[artifact, left=of test] (code) {\code{src/**}, \code{tests/test\_<group>.cpp},\\ \code{stress/stress\_main.cpp};\\ the suites on Windows + Linux,\\ the audits (coverage, KLM, consistency)};
  \node[agent, left=of code] (review) {Review\\ \footnotesize \code{code-reviewer} (read-only):\\ \footnotesize findings with failure scenarios};
  \node[artifact, below=8mm of review] (docs) {\code{CHANGELOG.md}, \code{README.md},\\ \code{ARCHITECTURE.md}, \code{USABILITY.md},\\ \code{MANUAL.md}, a Learn lesson,\\ the spec marked done};
  \node[artifact, right=of docs] (pr) {A branch \code{task/NNNN-slug}\\ (a subtree commit of \code{blendity/}),\\ a pull request on \code{main}\\ (\code{.github/pull\_request\_template.md})};
  \node[box, right=of pr] (ci) {CI: \code{build.yml}\\ \footnotesize windows, linux, macos,\\ \footnotesize linux-sanitizers};
  \node[gate, right=of ci] (g2) {The human\\ reviews and\\ merges};
  \node[box, below=8mm of g2] (main) {\code{main}\\ \footnotesize a merge commit;\\ \footnotesize \code{release.yml} builds artifacts};
  \node[gate, left=of main] (g3) {The user\\ asks for a\\ release};
  \node[artifact, left=of g3] (tag) {\code{VERSION} bumped, a tag \code{vX.Y.Z};\\ \code{release.yml}: installers,\\ \code{SHA256SUMS.txt}, the two books;\\ \code{publish\_release.sh}};
  \draw[flow] (ask) -- (plan);
  \draw[flow] (plan) -- (spec);
  \draw[flow] (spec) -- (g1);
  \draw[flow] (g1) -- (impl);
  \draw[flow] (impl) -- (test);
  \draw[flow] (test) -- (code);
  \draw[flow] (code) -- (review);
  \draw[flow] (review) -- (docs);
  \draw[flow] (docs) -- (pr);
  \draw[flow] (pr) -- (ci);
  \draw[flow] (ci) -- (g2);
  \draw[flow] (g2) -- (main);
  \draw[flow] (main) -- (g3);
  \draw[flow] (g3) -- (tag);
  \draw[dflow] (review.north) to[bend right=18] node[lbl, pos=0.5] {changes needed: fixed, or answered in the task's record} (impl.west);
  \draw[dflow] (ci.north) to[bend right=10] node[lbl, pos=0.5] {red: fixed on the branch} (pr.north);
\end{tikzpicture}
}
  \caption{One task from the request to the release. Diamonds are the human's gates; the boxes with paths are the artifacts.}
  \label{fig:sdlc}
\end{figure}

\begin{enumerate}
  \item \textbf{Plan.} A round of requests becomes task specs \code{docs/tasks/NNNN-slug.md} from \code{TEMPLATE.md}:
    goal, scope (in / out), likely files, acceptance criteria that a test can check, the tests and stress sections to
    run, risks and edge cases, the documentation to update. An architectural choice (a subsystem, a data-model change, a
    dependency, a change to a convention, to how the project is built, tested or released) gets an ADR from the
    template in \code{docs/decisions/README.md}: status, date, requirements, context, decision stated so review can
    check it, consequences. Records are never rewritten; a later one supersedes by number. \textbf{Gate: the human
    approves the spec and the ADR.}
  \item \textbf{Implement and test.} The implementer or the main session writes the code; every behaviour change gets a
    unit test in the \code{tests/test\_<group>.cpp} of its subsystem, named for what the user sees, in the round's group
    at the end of the file; a measurement is a stress row, never an asserted time (\adr{0056}). The suites run on
    Windows and Linux, both builds; the audits run (\code{test\_coverage\_map.py --check}, \code{hci\_consistency\_audit.py
    --check}, \code{hci\_klm.py --check}, \code{unity\_meta.py --check}, \code{make\_books.py --check}); geometry changes
    run the stress sections and their counts may not rise unexplained.
  \item \textbf{Review.} The reviewer's findings are fixed in a second pass, each with a test, or answered in the task's
    ``Review'' section. Some rounds took three passes; the record says so.
  \item \textbf{Document.} Every change: a CHANGELOG entry (asked / added / changed / checked, with the numbers), the
    README's count and feature paragraph, an ARCHITECTURE row for a new subsystem, a USABILITY row for a user-facing
    problem solved, the manual, a Learn lesson with Try-it buttons and questions for every new feature, the spec marked
    done with its PR link.
  \item \textbf{Publish.} A subtree commit of \code{blendity/} is pushed to a branch \code{task/NNNN-slug} and a pull
    request opened on \code{main} with the template (\code{.github/pull\_request\_template.md}): the task, what changed,
    the tests with their counts, the review's verdict, the documentation boxes. Commits end with the model's
    \code{Co-Authored-By} trailer; pull requests say they were generated with Claude Code. A round's PR is stacked on the
    previous round's branch when the rounds depend on each other, and the body says which PRs a merge also merges.
  \item \textbf{CI.} \code{build.yml} builds Windows, Linux and macOS, runs the unit tests and a quick stress pass, takes
    a headless screenshot, and runs the tests under AddressSanitizer and UndefinedBehaviorSanitizer on Linux; the four
    jobs are required checks. \code{main} is protected: pull requests only, no force pushes.
  \item \textbf{Merge.} \textbf{Gate: the human reviews and merges} with a merge commit, never a squash or a rebase of a
    stack; when the user asks an agent to merge, only after CI is green on all jobs. \code{release.yml} builds every
    merge as artifacts to try.
  \item \textbf{Release.} \textbf{Gate: the user asks for one.} \code{VERSION} is bumped, a \code{v*} tag pushed, and
    \code{release.yml} builds the installers, checks them, attaches the books and publishes (chapter~\ref{ch:build}).
\end{enumerate}

\section{Traceability}
\begin{figure}[H]
  \centering
  \resizebox{\textwidth}{!}{% Requirements traceability: how an ID in requirements.md reaches a test, and how the tools check the chain.
\begin{tikzpicture}[node distance=5mm and 9mm]
  \node[artifact] (req) {\code{docs/requirements.md}\\ R-01 \ldots\ R-15 (product),\\ Q-01 \ldots\ Q-05 (quality),\\ C-01 \ldots\ C-05 (constraints)};
  \node[artifact, right=of req] (task) {\code{docs/tasks/NNNN-*.md}\\ Requirements: R-xx, Q-xx\\ Decisions: ADR NNNN\\ Acceptance criteria 1..n\\ PR: the link};
  \node[artifact, right=of task] (adr) {\code{docs/decisions/NNNN-*.md}\\ Status, Date, Requirements\\ Context / Decision / Consequences\\ (never rewritten; superseded by number)};
  \node[artifact, below=of task] (test) {\code{tests/test\_<group>.cpp}\\ \code{test("...", []\{ CHECK(...) \})}\\ named for what the user sees;\\ \code{tests\_roundNNN} per round};
  \node[artifact, right=of test] (stress) {\code{stress/stress\_main.cpp}\\ limits, problem counts,\\ faces on top of each other,\\ timings (never asserted: \adr{0056})};
  \node[artifact, left=of test] (log) {\code{CHANGELOG.md}\\ one entry a round: asked,\\ added, changed, checked};
  \node[box, below=10mm of test, fill=blCode] (map) {\code{tools/test\_coverage\_map.py --check}\\ \footnotesize every console command, keymap action, window,\\ \footnotesize button and Inspector field against the tests that drive it};
  \node[box, left=of map, fill=blCode] (audits) {\code{hci\_consistency\_audit.py --check}\\ \code{hci\_klm.py --check}\\ \code{unity\_meta.py --check}\\ \code{make\_books.py --check}};
  \node[box, right=of map, fill=blCode] (ci) {CI (\code{build.yml})\\ \footnotesize the suites on three systems,\\ \footnotesize sanitizers, a quick stress pass,\\ \footnotesize a headless screenshot};
  \draw[flow] (req) -- (task);
  \draw[flow, <->] (task) -- (adr);
  \draw[flow] (task) -- (test);
  \draw[flow] (task) -- (stress);
  \draw[flow] (task) -- (log);
  \draw[dflow] (test) -- (map);
  \draw[dflow] (test) -- (ci);
  \draw[dflow] (log) -- (audits);
  \draw[dflow] (adr.east) to[bend left=30] (stress.east);
\end{tikzpicture}
}
  \caption{From a requirement ID to the tests and the audits that hold it.}
  \label{fig:trace}
\end{figure}

A requirement's ID (R-, Q-, C-) is cited by the task spec and the ADR; the spec's acceptance criteria become the
tests; the CHANGELOG entry names the task, the tests and the numbers; the coverage map ties every command, action,
window, button and field to the tests that drive it; and the HCI tools recount their tables so the documents cannot
drift from the counts. Changing a requirement is a human decision recorded in an ADR.

\section{Engineering standards}
\code{CLAUDE.md} is the contract every session reads first: the project map; how to build and run; the testing rules
(real events for UI behaviour, \code{meshops} calls for geometry, a scratch project, a test that disagrees with the code
is evidence); the code standards (C++17 on three compilers, match the surrounding code, copy-on-write, operators
through \code{LastOp}, every tool reachable from the console, no new dependency outside what Blender bundles, an SPDX
header in every file, robust geometry); the documentation duties; the workflow and approval gates; git and publishing
rules (commit only when asked, never the user's \code{Assets/}, never push to \code{main}, never rewrite history).

\section{The HCI method in the lifecycle}
Usability is designed and measured, not assumed (\adr{0085}): each HCI addition names the heuristic or model that
motivates it; the baseline is taken before any fix so each round's effect is a diff in git; the keystroke-level
estimate for frequent actions went from 80.9~s to 54.65~s over the pass (\code{docs/hci/klm.md}), the failing
walkthrough steps from 24 to 7, wrong shortcut labels to 0; the comparison reports put a VRChat avatar, a VRChat world
and a short film made in Blendity beside Blender, Unity, Godot, Unreal and Adobe from their documentation, with the
findings ranked for the user to choose from; the 118-task checklist carries target times. The think-aloud study with
SUS is written (\code{docs/hci/study-protocol.md}) and not yet run; the record says so.

\section{AI disclosure}
\code{docs/AI\_DISCLOSURE.md} states, for anyone reading, using or building on the project: the direction is human; the
code is written with Claude from approved specs under the standards above; every commit says so; numbers are measured,
not guessed, and where something was not tried the record says so; there is no telemetry, no cloud, and no AI at run
time except the optional local inbetweening download; what Blendity makes is the user's. Chapter~\ref{ch:licensing}
sets out why that is compatible with the GPL.
